%macbook 
\loai{Ví dụ mẫu}
\begin{vidu}%[3P2-6.1B][Loai1]
Cho hai nguồn kết hợp $A$ và $B$ cách nhau $13\ cm$ trên mặt nước dao động theo phương thẳng đứng với phương trình $u_1=u_2=4\cos 4\pi t\ cm$. Sóng do hai nguồn này tạo ra có bước sóng $3,2\ cm$. 
\begin{itemize}
\item[1.] Số điểm trên mặt nước dao động với biên độ cực đại và số điểm không dao động trên đoạn nối hai nguồn lần lượt là .......... và ..........
\item[2.] Điểm $M$ trên mặt nước cách $A$ và $B$ các khoảng lần lượt là $9\ cm$ và $6\ cm$. Số cực đại trên $MA$ là ..........
\item[3.] Dựng hình vuông $ABCD$ trên mặt nước. Lấy $C'$ đối xứng với $C$ qua $B$. Số cực đại trên $CC'$ là ..........
\item[4.] $I$ là trung điểm của $AB$. Vẽ đường tròn tâm $I$ nằm trên mặt nước với bán kính $10\ cm$. Số điểm không dao động trên đường tròn này là .......... 
\item[5.] Khoảng cách giữa hai điểm xa nhất có biên độ $8\ cm$ trên đoạn thẳng nối hai nguồn là .......... 
\item[6.] Điểm $P$ nằm trên mặt nước cách $A$ một khoảng $L$ và $AP$ vuông góc với $AB$. 
\begin{itemize}
\item[a.] Giá trị L nhỏ nhất để tại A dao động với biên độ cực đại là ..........
\item[b.] Giá trị L lớn nhất để tại A dao động với biên độ cực đại là ..........
\item[c.] Giá trị L nhỏ nhất để tại A dao động với biên độ cực tiểu là ..........
\item[d.] Giá trị L lớn nhất để tại A dao động với biên độ cực tiểu là ..........
\end{itemize}
\item[7.]  $I$ là trung điểm $AB$. $Q$ là điểm trên mặt nước nằm trên đường trung trực của $AB$ cách $I$ một đoạn $5\ cm$. Gọi $(d)$ là đường thẳng qua $Q$ và song song với $AB$. 
\begin{itemize}
\item[a.] Điểm $F$ thuộc $(d)$ và gần $Q$ nhất, dao động với biên độ cực đại. Khoảng cách $FQ$ là ..........
\item[b.]  Điểm $H$ thuộc $(d)$ và xa $Q$ nhất, dao động với biên độ cực đại. Khoảng cách $HQ$ là ..........
\item[c.]  Điểm $H$ thuộc $(d)$ và gần $Q$ nhất, dao động với biên độ cực tiểu. Khoảng cách $HQ$ là ..........
\item[d.]  Điểm $H$ thuộc $(d)$ và xa $Q$ nhất, dao động với biên độ cực tiểu. Khoảng cách $HQ$ là ..........
\end{itemize}
\end{itemize}
\loigiai{
\begin{itemize}
\item[1.] Số cực đại trên đoạn AB: $$-\dfrac{\ell}{\lambda}<k<\dfrac{\ell}{\lambda}\Leftrightarrow -4,1<k<4,1$$ Vậy có 9 cực đại trên AB\\
Số cực tiểu trên AB: $$-\dfrac{\ell}{\lambda}-\dfrac{1}{2}<k<\dfrac{\ell}{\lambda}-\dfrac{1}{2}\Leftrightarrow -4,6<k<3,6$$ Vậy có 8 cực tiểu trên AB
\item[2.] Số cực đại trên MA: $$-\dfrac{\ell}{\lambda}<k\le\dfrac{MA-MB}{\lambda}\Leftrightarrow -4,1<k\le 0,9$$ Vậy có 5 cực đại trên AM
\item[3.] Số cực đại trên BC: $$\dfrac{CA-CB}{\lambda}\le k<\dfrac{\ell}{\lambda}\Leftrightarrow 1,6\le k<4,1$$ Vậy có 3 cực đại trên BC. Do đó có 6 cực đại trên CC'
\item[4.] Đường tròn tâm I bán kính 10 cm bao cả hai nguồn bên trong\\
Mỗi đường cực tiểu cắt đường tròn tại hai điểm nên số cực tiểu trên đường tròn là $8.2=16$ điểm
\item[5.]
\item[6.] 
\item[7.] 
\end{itemize}
}
\end{vidu} %\dotlineEX{37}
 
\loai{Tìm số cực đại, cực tiểu trên đoạn thẳng khi biết cụ thể vị trí hai điểm đầu, cuối}
\begin{vidu}%[3P2-6.1K][Loai1]
\nguon{QG - 2019} Ở mặt chất lỏng, tại hai điểm $S_1$ và $S_2$ có hai nguồn dao động cùng pha dao động theo phương thẳng đứng phát ra hai sóng kết hợp có bước sóng 1 cm. Trong miền giao thoa, M là điểm cách $S_1$ và $S_2$ lần lượt 8 cm và 12 cm. Giữa M và đường trung trực của đoạn thẳng $S_1S_2$ có số vân giao thoa cực tiểu là 
\choice
{3}
{5}
{6}
{\True 4}
\loigiai{
\begin{itemize}
\item Gọi O là trung điểm của $S_1S_2$.
\item Vì hai nguồn cùng pha nên vân giao thoa cực tiểu thỏa mãn $d_1-d_2=\left(k+0,5 \right)\lambda$. 
\item Số vân cực tiểu giữa M và O là số giá trị k thõa: \begin{align*}
S_2O-S_1O<\left( k+0,5 \right)\lambda <S_2M-S_1M \Leftrightarrow 0<k+0,5<\dfrac{12-8}{1}\Rightarrow -0,5<k<3,5\Rightarrow k=0,\,1,\,2,\,3.
\end{align*}
\item Vậy có 4 đường cực tiểu thõa mãn.

\end{itemize}
}
\end{vidu} %\dotlineEX{5}
\loai{Tìm số cực đại, cực tiểu trên cạnh hình vuông}
\begin{vidu}%[3P2-6.1K][Loai1]
Thực hiện thí nghiệm giao thoa trên mặt nước với hai nguồn kết hợp A, B có khoảng cách $AB=8$ cm, phương trình sóng tại A, B là $x_A=x_B=a\cos40\pi t$ cm, vận tốc truyền sóng trên mặt nước là $v=30$ cm/s. Gọi C, D là hai điểm trên mặt nước sao cho ABCD là hình vuông. Số điểm dao động với biên độ cực đại trên CD là
\choice
{\True 5 điểm}
{11 điểm}
{10 điểm}
{7 điểm}
\loigiai{
\begin{itemize}
\item Ta có: $\lambda  = \dfrac{v}{f} = \dfrac{30}{20} = 1,5 $ cm.
\item Áp dụng định lí Pitago $AC = BD = 8\sqrt{2} $.
\item Số điểm cực đại trên CD thoả mãn:
\begin{align*}
\dfrac{\Delta d_C}{\lambda } \le k \le \dfrac{\Delta d_D}{\lambda} \Rightarrow \dfrac{CB-CA}{\lambda } \le k \le \dfrac{DB-DA}{\lambda} \Rightarrow \dfrac{8-8\sqrt{2}}{1,5} \le k \le \dfrac{8\sqrt{2}-8}{1,5} \Rightarrow -2,2 \le k \le 2,2.
\end{align*}
\item Vậy có 5 điểm dao động với biên độ cực đại trên CD.
\end{itemize}
}
\end{vidu} %\dotlineEX{5}
\loai{Tìm số cực đại, cực tiểu trên đoạn thẳng có 1 đầu thuộc đường trung trực}
\begin{vidu}%[3P2-6.1K][Loai1]
\nguon{MH - 2019} Trên mặt nước, tại hai điểm A và B cách nhau 19 cm, có hai nguồn kết hợp dao động cùng pha theo phương thẳng đứng, phát ra hai sóng có bước sóng 4 cm. Trong vùng giao thoa, M là một điểm ở mặt nước thuộc đường trung trực của AB. Trên đoạn AM, số điểm cực tiểu giao thoa là
\choice
{7}
{4}
{\True 5}
{6}
\loigiai{
\immini{\begin{itemize}
\item Số điểm dao động với biên độ cực tiểu trên AB là:
\begin{align*}
-\dfrac{AB}{\lambda}-\dfrac{1}{2}<k<\dfrac{AB}{\lambda}-\dfrac{1}{2}\Rightarrow -5,75<k<4,25\Rightarrow\ \ct{có 10 điểm}
\end{align*}
\item Như vậy trên AO sẽ có 5 vân hypebol cắt AO đồng nghĩa là trên AM có 5 điểm cực tiểu giao thoa.
\item 
\item 
\item 
\end{itemize}}{\begin{tikzpicture}[scale=1,thick,>=stealth']
\tkzDefPoints{0/0/o, -2/0/s1, 2/0/s2, 0/1.5/m, 0/2.5/y, 0/-0.5/y'}
\tkzDrawSegments(s1,s2 y,y' s1,m)
\tkzDrawPoints[size=3,fill=white](s1,s2,m)
\node at (s1)[below]{$A$};
\node at (s2)[below]{$B$};
\node at (m)[right]{$M$};
\node at (o)[below right]{$O$};
\end{tikzpicture}}
}
\end{vidu} %\dotlineEX{7}
\loai{Tìm số cực đại, cực tiểu trên đường vuông góc với đoạn hai nguồn tại nguồn}
\begin{vidu}%[3P2-6.1K][Loai1]
Ở mặt thoáng của một chất lỏng có hai nguồn sóng kết hợp A và B cách nhau 20 cm, dao động theo phương thẳng đứng với phương trình $u_A = u_B = 2\cos(40\pi t )$ ($u_A$ và $u_B$ tính bằng mm, t tính bằng s). Biết tốc độ truyền sóng trên mặt chất lỏng là 30 cm/s. Số điểm dao động với biên độ cực đại, cực tiểu trên đường thẳng vuông góc với AB tại A thuộc mặt phẳng giao thoa lần lượt là
\choice
{\True 26; 26}
{26; 24}
{13; 13}
{24; 26}
\loigiai{
\begin{itemize}
\item Số dãy cực đại về một phía đường trung trực là $\left[ \dfrac{OA}{0,5\lambda } \right]\overset{hay}=\left[ \dfrac{AB}{\lambda } \right]=13$\\
$\Rightarrow$  Mỗi dãy cắt đường d (vuông góc AB tại A) tại 2 điểm\\
$\Rightarrow$  Có 26 điểm dao động với biên độ cực đại trên d.
\item Số dãy cực tiểu về một phía đường trung trực là $\left[ \dfrac{AB}{\lambda }+0,5 \right]=13$\\
$\Rightarrow$  Có 26 điểm dao động với biên độ cực tiểu trên d.
 
\end{itemize}
}
\end{vidu} %\dotlineEX{5}
\loai{Tìm số cực đại, cực tiểu trên đường tròn bao cả hai nguồn bên trong}
\begin{vidu}%[3P2-6.1K][Loai1]
Trong thí nghiệm giao thoa sóng trên mặt chất lỏng, hai nguồn kết hợp $S_1$, $S_2$  cùng pha, cùng biên độ, cách nhau 9,5 cm. Khoảng cách gần nhất giữa vị trí cân bằng của hai phần tử trên mặt nước dao động với biên độ cực đại thuộc đoạn nối $S_1$, $S_2$  là 1 cm. Trên mặt nước vẽ một đường tròn sao cho vị trí nguồn $S_1$, $S_2$ ở trong đường tròn đó. Trên đường tròn ấy có bao nhiêu điểm có biên độ dao động cực đại?
\choice
{20}
{9}
{\True 18}
{10}
\loigiai{
\begin{itemize}
\item Khoảng cách gần nhất giữa hai điểm cực đại trên đoạn $S_1S_2$ là $0,5\lambda =1\ cm\Rightarrow \lambda =2\ cm$.
\item Số dãy cực đại giao thoa $-\dfrac{{S_1S_2}}{\lambda }< k< \dfrac{{S_1S_2}}{\lambda }\Leftrightarrow -4,75< k< 4,75\Rightarrow $ có 9 dãy.
\item Mỗi dãy cực đại cắt đường tròn tại 2 điểm $\Rightarrow $ có 18 điểm cực đại trên đường tròn.

\end{itemize}
}
\end{vidu} %\dotlineEX{5}
\loai{Biện luận số cực đại, cực tiểu}






















%macbook 
\loai{Vị trí cực đại trên đoạn hai nguồn}
\begin{vidu}%[3P2-6.2K][Loai1]
Hai nguồn sóng kết hợp A và B trên mặt nước cách nhau 20 cm dao động theo phương thẳng đứng với phương trình $u_1=u_2=4\cos 4\pi t$ cm. Điểm thuộc đoạn AB cách trung điểm của AB đoạn gần nhất 1,5 cm luôn không dao động. Khoảng cách giữa hai điểm xa nhất có biên độ 8 cm trên đoạn thẳng nối hai nguồn bằng
\choice
{12,5 cm}
{\True 18 cm}
{18,5 cm}
{19 cm}
\loigiai{
\begin{itemize}
\item Hai nguồn kết hợp cùng pha, trung điểm của AB là một cực đại, khoảng cách từ cực đại này đến cực tiểu gần nhất là $\dfrac{\lambda}{4}$, hay $\dfrac{\lambda}{4}=1,5$ cm $\Rightarrow \lambda =6$ cm.
\item Các điểm trên AB có biên độ 8 cm chính là các cực đại.
\item Số cực đại: $-\dfrac{AB}{\lambda}<k<\dfrac{AB}{\lambda}\Rightarrow -3,3<k<3,3\Rightarrow k=-3,...,3$.
\item Từ cực đại ứng với $k=-3$ đến cực đại ứng với $k=+3$ có 6 khoảng $\dfrac{\lambda}{2}$ nên khoảng cách giữa hai cực đại đó là $6\dfrac{\lambda}{2}=18$ cm.
\end{itemize}
}
\end{vidu} %\dotlineEX{5}
\loai{Vị trí cực đại gần nhất trên đường vuông góc với đoạn hai nguồn tại nguồn}
\begin{vidu}%[3P2-6.2K][Loai1]
Hai nguồn sóng kết hợp, cùng pha $S_1$ và $S_2$ cách nhau 2,2 m phát ra hai sóng có bước sóng 0,4 m, một điểm A nằm trên mặt chất lỏng cách $S_1$ một đoạn L và $AS_1$ vuông góc với $S_1S_2$. Giá trị L nhỏ nhất để tại A dao động với biên độ cực đại là
\choice
{0,4 m}
{\True 0,21 m}
{5,85 m}
{0,1 m}
\loigiai{
\begin{itemize}
\item Số điểm dao động với biên độ cực đại trên $S_1S_2$thoả mãn $\dfrac{-S_1S2}{\lambda}< k < \dfrac{-S_1S_2}{\lambda} \Rightarrow -5,5 < k < 5,5$.
\item A thuộc đường vuông góc với $S_1S_2$ qua $S_1$, để A gần $S_1$ nhất dao động với biên độ cực đại thì A thuộc vân cực đại bậc $-5$.
\item Ta có
$\begin{cases} AS_1-AS_2 = -5.\lambda\\AS^2_1 + S_1S_2^2 = AS^2_2\end{cases}$
$\Leftrightarrow L^2 + S_1S_2^2 = (L + 2)^2\Rightarrow L = 0,21\ m$.
\item 
\item 
\item 
\end{itemize}
}
\end{vidu} %\dotlineEX{5}
\loai{Vị trí cực tiểu trên đường vuông góc với đoạn hai nguồn tại nguồn}
\loai{Vị trí cực đại trên đường song song với đoạn hai nguồn}
\begin{vidu}%[3P2-6.2G][Loai1]
Ở mặt nước có hai nguồn sóng A, B dao động theo phương vuông góc với mặt nước, có phương trình $u=a\cos\omega t$, cách nhau 20 cm với bước sóng 5 cm. I là trung điểm AB. P là điểm nằm trên mặt nước, thuộc đường trung trực của AB và cách I một đoạn 5 cm. Gọi (d) là đường thẳng qua P và song song với AB. Điểm M thuộc (d) và gần P nhất, dao động với biên độ cực đại. Khoảng cách MP là
\choice
{2,5 cm}
{\True 2,81 cm}
{3 cm}
{3,81 cm}
\loigiai{
\immini{\begin{itemize}
\item Ta có: Để M dao động với biên độ cực đại và gần P nhất thì M là cực đại bậc 1
$d_1-d_2 = 1.\lambda = 5\quad (1)$
\item Mà $MQ^2 = d_1^2-AQ^2 = d_2^2-BQ^2$. Đặt $MP=QI=x$, ta có: $d_1^2-(10 + x)^2 = d_2^2-(10-x)^2\Rightarrow d_1 + d_2 = 8x\quad (2)$
\item Từ (1) và (2) $\Rightarrow 2d_1 = 5 + +8x \Rightarrow d_1 = 2,5 + 4x$.
\item Mặt khác: $d_1^2-(10 +x)^2 = 25\Rightarrow (2,5 + 4x)^2-(10 + x)^2 = 25\Rightarrow x = 2,81\ cm$.
\item 
\item 
\item 
\item 
\end{itemize}}{\begin{tikzpicture}[scale=1,thick,>=stealth']
\tkzDefPoints{0/0/i, 0/-0.5/y', 0/2.5/y, -3/0/a, 3/0/b, -3/2/c, 3/2/d, 0/2/p, 1.5/0/q, 1.5/2/m}
\tkzDrawSegments(y,y' a,b c,d a,m b,m m,q)
\node at (a)[below]{$A$};
\node at (b)[below]{$B$};
\node at (p)[above right]{$P$};
\node at (m)[above]{$M$};
\node at (q)[below]{$Q$};
\node at (i)[below right]{$I$};
\end{tikzpicture}}
}
\end{vidu} %\dotlineEX{5}
\loai{Vị trí cực đại trên đường tròn đường kính là đoạn hai nguồn}
\begin{vidu}%[3P2-6.2G][Loai1]
Trong thí nghiệm giao thoa sóng trên mặt nước, hai nguồn giống hệt nhau A và B cách nhau 7 cm, tạo ra sóng trên mặt nước với bước sóng 2 cm. Điểm M trên đường tròn đường kính AB (không nằm trên trung trực của AB) thuộc mặt nước xa đường trung trực của AB nhất dao động với biên độ cực đại. M cách A một đoạn nhỏ nhất, lớn nhất lần lượt là
\choice
{0,94 cm và 6,57 cm}
{\True 0,94 cm và 6,94 cm}
{0,87 cm và 6,94 cm}
{1,77 cm và 6,77 cm}
\loigiai{
\begin{itemize}
\item Hai cực đại xa nhất nằm hai bên đường trung trực có hiệu đường đi $MA-MB=-n\lambda$ (M gần A hơn) và $MA-MB=n\lambda$ (M xa A hơn); với n là số nguyên lớn nhất thỏa mãn $n<\dfrac{AB}{\lambda}=\dfrac{7}{2}=3,5\Rightarrow n=3$.
\item $\begin{cases}
MA-\sqrt{AB^2-MA^2}=-3\lambda \Rightarrow MA-\sqrt{7^2-MA^2}=-6\Rightarrow MA=0,94\ cm.\\
MA-\sqrt{AB^2-MA^2}=3\lambda \Rightarrow MA-\sqrt{7^2-MA^2}=6\Rightarrow MA=6,94\ cm.
\end{cases}$

\end{itemize}
}
\end{vidu} %\dotlineEX{5}
\loai{Vị trí cực tiểu trên đường tròn đường kính là đoạn hai nguồn}
\begin{vidu}%[3P2-6.2K][Loai1]
Trong hiện tượng giao thoa sóng nước, hai nguồn dao động theo phương vuông góc với mặt nước, cùng biên độ, cùng pha, cùng tần số 15 Hz được đặt tại hai điểm $S_1$ và $S_2$ cách nhau 10 cm. Xét các điểm trên mặt nước thuộc đường tròn đường kính $S_1S_2$, điểm mà phần tử tại đó dao động với biên độ cực tiểu cách đường trung trực của  $S_1S_2$ một đoạn ngắn nhất là 1,4 cm. Tốc độ truyền sóng trên bề mặt chất lỏng là
\choice
{0,42 m/s}
{\True 0,6 m/s}
{0,3 m/s}
{0,84 m/s}
\loigiai{
\immini{\begin{itemize}
\item Rõ ràng M thuộc đường cực tiểu số 1 $\Rightarrow MS_1 - MS_2 = 0,5.\lambda\quad (*)$
\item Bài cho $MI = OH = 1,4$ cm mà $\Delta MS_1S_2$ vuông tại M\\
$\Rightarrow M{{S}_{1}}=\sqrt{S_1H.S_1S_2}=$ 8 cm và $M{{S}_{2}}=\sqrt{S_2H.S_1S_2}=$ 6 cm.
\item Từ $(*) \Rightarrow  \lambda  = 4\ cm \Rightarrow v = 60\ cm/s$.
\item 
\item 
\item 
\end{itemize}}{\begin{tikzpicture}[scale=0.6,thick,>=stealth']
\tkzInit[ymin=-3.5,ymax=5,xmin=-1,xmax=7]\tkzClip
\tkzDefPoints{0/0/a, 3/0/o, 6/0/b, 3/4/p, 3/-4/q, 5/4/y, 5/-4/y'}
\tkzDefShiftPoint[o](63:3){m}
\tkzDefPointBy[projection=onto a--b](m) \tkzGetPoint{h}
\tkzDefPointBy[projection=onto p--q](m) \tkzGetPoint{i}
\tkzDrawCircle[radius](o,a)
\tkzDrawSegments[dashed](a,b a,m b,m h,m i,m)
\tkzDrawSegments[dashed](p,q)
\tkzDrawPoints[size=3,fill=white](a,b,o,m,i,h)
\tkzMarkRightAngles(a,m,b)
\draw[blue](y)node[above,black]{$k=1$} to [bend right =30](y');
\node at (a)[below left]{$S_1$};
\node at (b)[below right]{$S_2$};
\node at (m)[right,yshift=3pt]{$M$};
\node at (o)[below left]{$O$};
\node at (h)[below]{$H$};
\node at (i)[left]{$I$};
\end{tikzpicture}}
}
\end{vidu} %\dotlineEX{5}
\loai{Vị trí cực đại xa (gần) đường thẳng chứa hai nguồn nhất trên đường tròn bán kính là đoạn hai nguồn}
\begin{vidu}%[3P2-6.2K][Loai1]
Trong hiện tượng giao thoa sóng nước, hai nguồn A, B cách nhau 20 cm dao động cùng biên độ, cùng pha, cùng tần số 50 Hz. Tốc độ truyền sóng trên mặt nước là 1,5 m/s. Xét các điểm trên mặt nước thuộc đường tròn tâm A, bán kính AB, điểm dao động với biên độ cực đại cách đường thẳng AB một đoạn ngắn nhất bằng
\choice
{18,67 mm}
{17,96 mm}
{\True 19,97 mm}
{15,39 mm}
\loigiai{
\immini{\begin{itemize}
\item $\lambda =\dfrac{v}{f}=3(cm)$
$n = \text{số nguyên lớn nhất}<\dfrac{AB}{\lambda}
=\dfrac{20}{3}=6,7\Rightarrow n=6$.
\item Điểm M phải là cực đại gần B nhất nên:
$MA-MB=6\lambda \Rightarrow MB=2\ cm \cos \alpha =\dfrac{AB^2+MB^2-MA^2}{2.AB.MB}=0,995;\ NH=AN.\sin \alpha =AN.\sqrt{1-{\cos}^2\alpha}=20.\sqrt{1-0,995^2}=1,997\ cm$.
\end{itemize}}{\begin{tikzpicture}[scale=1,thick,>=stealth']
\tkzInit[ymin=-3,ymax=3,xmin=-3,xmax=3]\tkzClip
\tkzDefPoints{0/0/s1, 1/0/i, 2/0/s2, 1/3/y, 1/-3/y', 0/3/a, 0/-3/b, 2/3/c, 2/-3/d}
\tkzDefShiftPoint[s1](40:2){m}
\tkzDefPointBy[projection=onto y--y'](m) \tkzGetPoint{j}
\tkzDefPointBy[projection=onto s1--s2](m) \tkzGetPoint{h}
\tkzDefShiftPoint[i](0:0.4){n}
\tkzDrawSegments(y,y' s1,s2 s1,m s2,m h,m)
\tkzDrawCircle[radius](s1,s2)
\tkzDrawSegments[very thick,red](m,j)
\tkzDrawPoints[size=3](s1,s2,i,m,j,h)
\draw[blue](c) to [bend right =20](d);
\tkzMarkAngles[size=0.4cm](s2,s1,m)
\tkzLabelAngle[pos=0.7](s2,s1,m){$\alpha$}
\node at (s1)[below left]{$A$};
\node at (s2)[below right]{$B$};
\node at (i)[below left]{$O$};
\node at (h)[below]{$H$};
\node at (m)[right]{$M$};
\node at (j)[left]{$J$};
\end{tikzpicture}}
}
\end{vidu} %\dotlineEX{5}

\loai{Biện luận tìm các giá trị}
\begin{vidu}%[3P2-6.2G][Loai1]
Ở mặt chất lỏng có hai nguồn sóng A, B cách nhau 24 cm, dao động theo phương thẳng đứng với phương trình là $u_A = u_B = a\cos60\pi t$ (với t tính bằng s). Tốc độ truyền sóng trên mặt chất lỏng là 45 cm/s. Gọi $MN = 4$ cm là đoạn thẳng trên mặt chất lỏng có chung trung trực với AB. Khoảng cách xa nhất giữa MN với AB là bao nhiêu để có ít nhất 5 điểm dao động cực đại nằm trên MN?
\choice
{12,7 cm}
{\True 10,5 cm}
{14,2 cm}
{6,4 cm}
\loigiai{
\immini{\begin{itemize}
\item Ta thấy  M và N phải thuộc dãy cực đại số 2 (tính từ đường trung trực\\
$\Rightarrow MA - MB = 2.\lambda\  (*) \Rightarrow \sqrt{x^2+14^2}-\sqrt{x^2+10^2}= 3\ cm\Rightarrow  x = 10,5\ cm$.\\
Hoặc: $MA^2 - MB^2 = AH^2 - HB^2 = 96\ cm$.
\item Từ $(*) \Rightarrow  MA + MB = 32\ cm \Rightarrow  MB = 14,5\ cm \Rightarrow  x = MH = 10,5\ cm$.
\end{itemize}}{\begin{tikzpicture}[scale=1]
\begin{scope}[shift={(0,0)}]
\tkzInit[ymin=-3,ymax=3,xmin=-3,xmax=3]\tkzClip
\tkzDefPoints{-2/0/s1, 2/0/s2,1.32/1.5/m,-1.32/1.5/n}
\tkzDefPointBy[projection=onto s1--s2](m) \tkzGetPoint{h}
\draw[red] plot[domain=-2:2] ({1*cosh(\x)-0.5},{3*sinh(\x)});
\draw[red] plot[domain=-2:2] ({-1*cosh(\x)+0.5},{3*sinh(\x)});
\draw[red] plot[domain=-2:2] ({1*cosh(\x)},{1.75*sinh(\x)});
\draw[red] plot[domain=-2:2] ({-1*cosh(\x)},{1.75*sinh(\x)});
\draw[red] plot[domain=-2:2] ({1*cosh(\x)+0.5},{1.01*sinh(\x)});
\draw[red] plot[domain=-2:2] ({-1*cosh(\x)-0.5},{1.01*sinh(\x)});
\draw[red](0,3)--(0,-3);
\draw[blue](s1)--(s2)(m)--(n);
\draw[blue,dashed](s1)--(m)(s2)--(m)(h)--node[midway,left,black]{$x$}(m);
\tkzDrawPoints[size=3,fill=white](s1,s2,m,n,h)
\tkzLabelPoint[left](s1){$A$}
\tkzLabelPoint[right](s2){$B$}
\tkzLabelPoint[above left](m){$M$}
\tkzLabelPoint[above right](n){$N$}
\tkzLabelPoint[below](h){$H$}
\end{scope}
\end{tikzpicture}}
}
\end{vidu} %\dotlineEX{5}































